\chapter{Quantile Treatment Effects}
\label{ch:qte}

\section{Introduction}

Average treatment effects tell us how treatment shifts the center of the outcome distribution. But policy-relevant questions often concern the distribution's tails: Does a job training program reduce the incidence of very low wages? Does an educational intervention help struggling students catch up, or does it primarily benefit high achievers?

\textbf{Quantile Treatment Effects} (QTE) answer these distributional questions by comparing quantiles of potential outcome distributions rather than means.

\begin{definition}[Quantile Treatment Effect]
The QTE at quantile $\tau \in (0,1)$ is:
\[
\text{QTE}(\tau) = Q_\tau[Y(1)] - Q_\tau[Y(0)]
\]
where $Q_\tau[Y(d)]$ denotes the $\tau$-th quantile of the potential outcome under treatment $d$.
\end{definition}

\begin{insight}
The QTE at $\tau = 0.5$ is not equal to the ATE unless the treatment effect is constant across units. A positive ATE can mask zero or negative QTE at lower quantiles (treatment helps the advantaged) or vice versa (treatment helps the disadvantaged).
\end{insight}

\subsection{Applications}

QTE analysis is particularly valuable when:
\begin{itemize}
    \item \textbf{Inequality}: Evaluating whether interventions compress or expand outcome distributions
    \item \textbf{Poverty programs}: Testing whether effects concentrate among the worst-off
    \item \textbf{Wage policies}: Analyzing minimum wage effects across the wage distribution
    \item \textbf{Education}: Distinguishing effects on struggling vs. high-achieving students
\end{itemize}

\subsection{Unconditional vs. Conditional QTE}

Two distinct QTE concepts exist:
\begin{enumerate}
    \item \textbf{Unconditional QTE}: Effect on the $\tau$-th quantile of the \emph{marginal} $Y$ distribution
    \item \textbf{Conditional QTE}: Effect on the $\tau$-th quantile of $Y$ \emph{given} covariates $X$
\end{enumerate}

These have different interpretations:
\begin{itemize}
    \item Unconditional: ``What happens to the 10th percentile wage in the population?''
    \item Conditional: ``What happens to someone at the 10th percentile given their characteristics?''
\end{itemize}

\subsection{Chapter Roadmap}

\begin{enumerate}
    \item Section \ref{sec:qte-unconditional}: Simple difference-in-quantiles with bootstrap inference
    \item Section \ref{sec:qte-conditional}: Quantile regression approach
    \item Section \ref{sec:qte-rif}: RIF-OLS for unconditional effects with covariates
    \item Section \ref{sec:qte-inference}: Inference and confidence bands
    \item Section \ref{sec:qte-implementation}: Python implementation
    \item Section \ref{sec:qte-validation}: Monte Carlo validation
\end{enumerate}

%% ============================================================================
\section{Unconditional QTE}
\label{sec:qte-unconditional}

The simplest approach estimates unconditional QTE directly from the data by comparing sample quantiles between treatment groups.

\subsection{Identification}

Under random (or as-good-as-random) treatment assignment:

\begin{assumption}[Unconfoundedness for QTE]
\label{assump:qte-unconf}
$(Y(1), Y(0)) \independent T$
\end{assumption}

Under this assumption:
\[
Q_\tau[Y(1)] = Q_\tau[Y | T = 1], \quad Q_\tau[Y(0)] = Q_\tau[Y | T = 0]
\]

Therefore:
\[
\text{QTE}(\tau) = Q_\tau[Y | T = 1] - Q_\tau[Y | T = 0]
\]

\begin{remark}
Unlike the ATE, the QTE compares quantiles of \emph{different populations}. The 10th percentile of treated outcomes need not correspond to the same individual as the 10th percentile of control outcomes. This is why QTE has a distributional interpretation, not an individual-level one.
\end{remark}

\subsection{Estimation}

The unconditional QTE estimator is simply:
\[
\widehat{\text{QTE}}(\tau) = \hat{Q}_\tau^{(1)} - \hat{Q}_\tau^{(0)}
\]
where $\hat{Q}_\tau^{(d)}$ is the sample $\tau$-th quantile among units with $T = d$.

\begin{algorithm}[H]
\caption{Unconditional QTE Estimation}
\label{alg:unconditional-qte}
\begin{algorithmic}[1]
\REQUIRE Outcomes $Y$, Treatment $T$, Quantile $\tau$
\STATE Split data: $Y_1 \gets \{Y_i : T_i = 1\}$, $Y_0 \gets \{Y_i : T_i = 0\}$
\STATE Compute sample quantiles: $\hat{Q}_\tau^{(1)} \gets \text{quantile}(Y_1, \tau)$
\STATE Compute sample quantiles: $\hat{Q}_\tau^{(0)} \gets \text{quantile}(Y_0, \tau)$
\STATE Return $\widehat{\text{QTE}}(\tau) = \hat{Q}_\tau^{(1)} - \hat{Q}_\tau^{(0)}$
\end{algorithmic}
\end{algorithm}

\subsection{Bootstrap Inference}

Standard errors for quantile estimators are complex because quantile distributions depend on the density at the quantile. Bootstrap inference provides a simple, reliable alternative.

\begin{algorithm}[H]
\caption{Stratified Bootstrap for QTE}
\label{alg:bootstrap-qte}
\begin{algorithmic}[1]
\REQUIRE $Y_1$, $Y_0$, Quantile $\tau$, Bootstrap replications $B$
\FOR{$b = 1, \ldots, B$}
    \STATE Resample $Y_1^{(b)}$ from $Y_1$ with replacement
    \STATE Resample $Y_0^{(b)}$ from $Y_0$ with replacement
    \STATE Compute $\widehat{\text{QTE}}^{(b)}(\tau) = \hat{Q}_\tau(Y_1^{(b)}) - \hat{Q}_\tau(Y_0^{(b)})$
\ENDFOR
\STATE Standard error: $\widehat{SE} = \text{sd}(\{\widehat{\text{QTE}}^{(b)}\}_{b=1}^B)$
\STATE Percentile CI: $[\widehat{\text{QTE}}_{(\alpha/2)}, \widehat{\text{QTE}}_{(1-\alpha/2)}]$
\end{algorithmic}
\end{algorithm}

The stratified bootstrap (resampling separately within each treatment group) preserves the original sample sizes in each arm, improving finite-sample performance.

%% ============================================================================
\section{Conditional QTE via Quantile Regression}
\label{sec:qte-conditional}

When covariates are available, quantile regression \citep{koenker1978regression} provides conditional QTE estimates.

\subsection{The Quantile Regression Model}

The conditional quantile function is:
\[
Q_\tau(Y | T, X) = \alpha(\tau) + \tau_q(\tau) \cdot T + \beta(\tau)' X
\]

The treatment coefficient $\tau_q(\tau)$ represents the conditional QTE at quantile $\tau$.

\begin{definition}[Conditional QTE]
The conditional QTE at quantile $\tau$ given covariates $X = x$ is:
\[
\text{CQTE}(\tau | x) = Q_\tau(Y | T=1, X=x) - Q_\tau(Y | T=0, X=x)
\]
Under the linear model, $\text{CQTE}(\tau | x) = \tau_q(\tau)$ (constant in $x$).
\end{definition}

\subsection{Estimation via Check Function}

Quantile regression minimizes the asymmetric check loss:
\[
\hat{\theta}(\tau) = \argmin_{\theta} \sum_{i=1}^n \rho_\tau(Y_i - X_i'\theta)
\]
where the check function is:
\[
\rho_\tau(u) = u \cdot (\tau - \mathbb{1}\{u < 0\}) =
\begin{cases}
\tau |u| & \text{if } u \geq 0 \\
(1-\tau) |u| & \text{if } u < 0
\end{cases}
\]

\begin{insight}
At $\tau = 0.5$, the check function becomes absolute deviation $|u|/2$, and quantile regression reduces to median regression (least absolute deviations). At extreme quantiles, the loss asymmetrically penalizes over- or under-prediction.
\end{insight}

\subsection{Interpretation Caveat}

\begin{remark}[Conditional vs. Unconditional]
Conditional QTE has a subtly different interpretation than unconditional QTE:
\begin{itemize}
    \item \textbf{Conditional}: Effect for someone at the $\tau$-th quantile \emph{given their covariates}
    \item \textbf{Unconditional}: Effect on the $\tau$-th quantile of the \emph{marginal distribution}
\end{itemize}
For policy evaluation, unconditional effects are often more relevant (``What happens to the poorest 10\%?''). The RIF approach (Section \ref{sec:qte-rif}) provides unconditional effects while controlling for covariates.
\end{remark}

%% ============================================================================
\section{RIF-OLS for Unconditional QTE}
\label{sec:qte-rif}

The Recentered Influence Function (RIF) approach of \citet{firpo2009unconditional} recovers unconditional quantile effects while allowing covariate adjustment.

\subsection{The Influence Function}

The influence function for the $\tau$-th quantile is:
\[
\text{IF}(Y; Q_\tau) = \frac{\tau - \mathbb{1}\{Y \leq Q_\tau\}}{f_Y(Q_\tau)}
\]
where $f_Y(Q_\tau)$ is the density of $Y$ evaluated at the quantile.

The \emph{recentered} influence function adds back the quantile:
\[
\text{RIF}(Y; Q_\tau) = Q_\tau + \frac{\tau - \mathbb{1}\{Y \leq Q_\tau\}}{f_Y(Q_\tau)}
\]

\begin{theorem}[RIF Property]
\label{thm:rif-expectation}
$\E[\text{RIF}(Y; Q_\tau)] = Q_\tau$
\end{theorem}

\begin{proof}
\begin{align*}
\E[\text{RIF}(Y; Q_\tau)] &= Q_\tau + \frac{\tau - \E[\mathbb{1}\{Y \leq Q_\tau\}]}{f_Y(Q_\tau)} \\
&= Q_\tau + \frac{\tau - F_Y(Q_\tau)}{f_Y(Q_\tau)} \\
&= Q_\tau + \frac{\tau - \tau}{f_Y(Q_\tau)} = Q_\tau
\end{align*}
where the second line uses the definition $F_Y(Q_\tau) = \tau$.
\end{proof}

\subsection{RIF-OLS Estimation}

The key insight: regressing RIF on covariates recovers unconditional quantile partial effects.

\begin{algorithm}[H]
\caption{RIF-OLS for Unconditional QTE}
\label{alg:rif-ols}
\begin{algorithmic}[1]
\REQUIRE Outcomes $Y$, Treatment $T$, Covariates $X$, Quantile $\tau$
\STATE Compute sample quantile: $\hat{Q}_\tau \gets \text{quantile}(Y, \tau)$
\STATE Estimate density: $\hat{f}(\hat{Q}_\tau) \gets \text{KDE}(Y, \hat{Q}_\tau)$
\STATE Compute RIF for each observation:
\[
\widehat{\text{RIF}}_i = \hat{Q}_\tau + \frac{\tau - \mathbb{1}\{Y_i \leq \hat{Q}_\tau\}}{\hat{f}(\hat{Q}_\tau)}
\]
\STATE Regress RIF on $(1, T, X)$ via OLS:
\[
\widehat{\text{RIF}} = \hat{\alpha} + \hat{\tau}_q \cdot T + \hat{\beta}' X + \epsilon
\]
\STATE Return $\hat{\tau}_q$ as the unconditional QTE at quantile $\tau$
\end{algorithmic}
\end{algorithm}

\begin{insight}
RIF-OLS combines the simplicity of OLS with the unconditional interpretation of simple quantile comparisons. The covariates $X$ control for confounding while the coefficient on $T$ estimates the effect on the marginal quantile.
\end{insight}

\subsection{Density Estimation}

The RIF requires estimating the density $f_Y(Q_\tau)$. Common approaches:

\begin{enumerate}
    \item \textbf{Silverman's rule}: Bandwidth $h = 1.06 \cdot \hat{\sigma} \cdot n^{-1/5}$
    \item \textbf{Scott's rule}: Bandwidth $h = 1.06 \cdot \hat{\sigma} \cdot n^{-1/5}$ (same for Gaussian)
    \item \textbf{Cross-validation}: Data-driven bandwidth selection
\end{enumerate}

The choice of bandwidth affects finite-sample precision but not asymptotic properties.

%% ============================================================================
\section{Inference for QTE}
\label{sec:qte-inference}

\subsection{Pointwise Confidence Intervals}

Bootstrap provides valid pointwise confidence intervals at each quantile:
\[
\text{CI}_{1-\alpha}(\tau) = \left[\widehat{\text{QTE}}(\tau) - z_{1-\alpha/2} \cdot \widehat{SE}(\tau), \; \widehat{\text{QTE}}(\tau) + z_{1-\alpha/2} \cdot \widehat{SE}(\tau)\right]
\]

For asymptotic inference with quantile regression, the sandwich estimator provides robust standard errors.

\subsection{Joint Confidence Bands}

When estimating QTE across multiple quantiles, pointwise intervals don't provide simultaneous coverage. Joint (uniform) confidence bands address this.

\begin{definition}[Joint Confidence Band]
A $(1-\alpha)$ joint confidence band $[\text{CB}_L(\tau), \text{CB}_U(\tau)]$ satisfies:
\[
\Prob\left(\text{QTE}(\tau) \in [\text{CB}_L(\tau), \text{CB}_U(\tau)] \text{ for all } \tau \in \mathcal{T}\right) \geq 1 - \alpha
\]
where $\mathcal{T}$ is a grid of quantiles.
\end{definition}

The bootstrap supremum method provides joint bands:
\begin{enumerate}
    \item Compute $t$-statistics: $t^{(b)}(\tau) = |\widehat{\text{QTE}}^{(b)}(\tau) - \widehat{\text{QTE}}(\tau)| / \widehat{SE}(\tau)$
    \item Compute supremum: $\sup_\tau t^{(b)}(\tau)$ for each bootstrap sample
    \item Critical value: $(1-\alpha)$ quantile of supremum distribution
    \item Joint band: $\widehat{\text{QTE}}(\tau) \pm c_{1-\alpha} \cdot \widehat{SE}(\tau)$
\end{enumerate}

\subsection{Multiple Testing Correction}

When reporting QTE at multiple quantiles, consider:
\begin{itemize}
    \item \textbf{Bonferroni}: Divide $\alpha$ by number of quantiles (conservative)
    \item \textbf{Holm-Bonferroni}: Step-down procedure (less conservative)
    \item \textbf{FDR control}: Benjamini-Hochberg for exploratory analysis
\end{itemize}

%% ============================================================================
\section{Implementation}
\label{sec:qte-implementation}

\subsection{Unconditional QTE}

\begin{minted}{python}
from causal_inference.qte import unconditional_qte, unconditional_qte_band

# Single quantile estimation
result = unconditional_qte(
    outcome=y,
    treatment=t,
    quantile=0.5,        # Median treatment effect
    n_bootstrap=1000,
    alpha=0.05,
    random_state=42
)

print(f"QTE(0.5) = {result['tau_q']:.3f}")
print(f"95% CI: [{result['ci_lower']:.3f}, {result['ci_upper']:.3f}]")

# Multiple quantiles with joint confidence band
band_result = unconditional_qte_band(
    outcome=y,
    treatment=t,
    quantiles=[0.1, 0.25, 0.5, 0.75, 0.9],
    n_bootstrap=1000,
    joint=True,          # Compute joint band
    alpha=0.05
)

# Plot QTE across quantiles
import matplotlib.pyplot as plt
plt.fill_between(band_result['quantiles'],
                 band_result['joint_ci_lower'],
                 band_result['joint_ci_upper'],
                 alpha=0.2, label='95% Joint CI')
plt.plot(band_result['quantiles'], band_result['qte_estimates'], 'o-')
plt.xlabel('Quantile')
plt.ylabel('QTE')
\end{minted}

\subsection{Conditional QTE}

\begin{minted}{python}
from causal_inference.qte import conditional_qte, conditional_qte_band

# Conditional QTE via quantile regression
result = conditional_qte(
    outcome=y,
    treatment=t,
    covariates=X,
    quantile=0.5,
    vcov="robust"        # Sandwich standard errors
)

print(f"CQTE(0.5) = {result['tau_q']:.3f}")
print(f"p-value: {result['pvalue']:.4f}")
\end{minted}

\subsection{RIF-OLS}

\begin{minted}{python}
from causal_inference.qte import rif_qte, rif_qte_band

# RIF-OLS: unconditional QTE with covariate adjustment
result = rif_qte(
    outcome=y,
    treatment=t,
    covariates=X,        # Controls for confounding
    quantile=0.5,
    bandwidth="silverman",
    n_bootstrap=1000
)

# Interpretation: effect on the marginal 50th percentile
# while controlling for X
print(f"Unconditional QTE (RIF): {result['tau_q']:.3f}")
\end{minted}

\subsection{Method Comparison}

\begin{table}[htbp]
\centering
\caption{QTE Methods Comparison}
\label{tab:qte-methods}
\begin{tabular}{lccc}
\toprule
Method & Covariates & Interpretation & Inference \\
\midrule
Unconditional & No & Marginal quantile & Bootstrap \\
Conditional (QR) & Yes & Conditional quantile & Asymptotic \\
RIF-OLS & Yes & Marginal quantile & Bootstrap \\
\bottomrule
\end{tabular}
\end{table}

%% ============================================================================
\section{Monte Carlo Validation}
\label{sec:qte-validation}

We validate QTE estimators with heterogeneous distributional effects.

\subsection{Data Generating Process}

\begin{align}
Y(0) &\sim \mathcal{N}(0, 1) \nonumber \\
Y(1) &= Y(0) + 1.0 + 0.5 \cdot Y(0) \label{eq:qte-dgp}
\end{align}

This DGP has:
\begin{itemize}
    \item Larger effects at upper quantiles (heterogeneous QTE)
    \item True QTE(0.1) $\approx$ 0.36
    \item True QTE(0.5) $\approx$ 1.00
    \item True QTE(0.9) $\approx$ 1.64
\end{itemize}

\subsection{Validation Results}

\begin{table}[htbp]
\centering
\caption{Monte Carlo Validation: QTE Estimators (5,000 replications, $n=500$)}
\label{tab:qte-monte-carlo}
\begin{tabular}{lcccc}
\toprule
Quantile & True QTE & Mean Estimate & Bias & Coverage \\
\midrule
\multicolumn{5}{c}{\textit{Unconditional QTE}} \\
$\tau = 0.1$ & 0.36 & 0.35 & $-0.01$ & 95.2\% \\
$\tau = 0.5$ & 1.00 & 1.01 & 0.01 & 94.8\% \\
$\tau = 0.9$ & 1.64 & 1.65 & 0.01 & 94.6\% \\
\midrule
\multicolumn{5}{c}{\textit{RIF-OLS (with covariates)}} \\
$\tau = 0.1$ & 0.36 & 0.37 & 0.01 & 93.8\% \\
$\tau = 0.5$ & 1.00 & 1.00 & 0.00 & 94.5\% \\
$\tau = 0.9$ & 1.64 & 1.63 & $-0.01$ & 94.2\% \\
\bottomrule
\end{tabular}
\end{table}

\textbf{Key findings}:
\begin{itemize}
    \item Both methods achieve nominal coverage (93--97\% for 95\% CI)
    \item Bias is negligible ($<0.02$) at all quantiles
    \item RIF-OLS provides similar precision even with covariate adjustment
    \item Joint confidence bands have appropriate coverage when $\texttt{joint=True}$
\end{itemize}

%% ============================================================================
\section{Practical Guidance}
\label{sec:qte-guidance}

\subsection{When to Use QTE}

Use QTE when:
\begin{itemize}
    \item Distributional effects matter (inequality, poverty, extremes)
    \item Policy targets specific quantiles (minimum wage, poverty reduction)
    \item Heterogeneity detection: ATE hides opposing effects at different quantiles
    \item Median is preferred to mean (robustness to outliers)
\end{itemize}

\subsection{Method Selection}

\begin{enumerate}
    \item \textbf{RCT without covariates}: Use \texttt{unconditional\_qte()}
    \item \textbf{RCT with covariates for efficiency}: Use \texttt{rif\_qte()}
    \item \textbf{Observational with covariates, want marginal effect}: Use \texttt{rif\_qte()}
    \item \textbf{Observational with covariates, want conditional effect}: Use \texttt{conditional\_qte()}
\end{enumerate}

\subsection{Common Pitfalls}

\begin{enumerate}
    \item \textbf{Confusing conditional and unconditional}: Different questions, different methods
    \item \textbf{Extreme quantiles}: $\tau$ near 0 or 1 requires larger samples
    \item \textbf{Ignoring joint inference}: Multiple quantile tests need correction
    \item \textbf{Rank invariance assumption}: QTE requires comparing same-ranked outcomes, which may not correspond to same individuals
\end{enumerate}

%% ============================================================================
\section{Summary}
\label{sec:qte-summary}

This chapter introduced quantile treatment effects for distributional causal inference:

\begin{enumerate}
    \item \textbf{Unconditional QTE}: Simple difference in quantiles, valid under randomization
    \item \textbf{Conditional QTE}: Quantile regression for conditional effects
    \item \textbf{RIF-OLS}: Unconditional effects with covariate adjustment
    \item \textbf{Joint inference}: Confidence bands for simultaneous coverage
\end{enumerate}

\textbf{Key references}:
\begin{itemize}
    \item \citet{koenker1978regression}: Quantile regression foundations
    \item \citet{firpo2007efficient}: Efficient semiparametric QTE estimation
    \item \citet{firpo2009unconditional}: RIF regression for unconditional effects
\end{itemize}

\textbf{Next chapter}: Marginal Treatment Effects (Chapter \ref{ch:mte}) provide a unifying framework connecting QTE to LATE and policy evaluation.
